\section{Módulo Python de Detección}

\subsection{vehiculo\_detector.py}

\textbf{Archivo}: \texttt{deteccion/vehiculo\_detector.py} (767 líneas)

\textbf{Responsabilidad}: Sistema completo de detección vehicular con YOLOv8, tracking multi-cámara, WebSocket streaming, y persistencia XML.

\subsubsection{Arquitectura del Detector}

\begin{figure}[H]
\centering
\begin{tikzpicture}[
    box/.style={rectangle, draw, fill=cyan!20, minimum width=2.5cm, minimum height=0.8cm, font=\small},
    node distance=1cm
]
    \node[box] (cam1) {CAM\_001};
    \node[box, right=of cam1] (cam2) {CAM\_002};
    \node[box, below=of cam1, xshift=1.25cm] (workers) {Workers};
    \node[box, below=of workers] (yolo) {YOLOv8};
    \node[box, below=of yolo, xshift=-1.5cm] (track) {Tracking};
    \node[box, below=of yolo, xshift=1.5cm] (color) {Color};
    \node[box, below=of yolo] (xml) {XML DB};
    
    \draw[-{Latex}] (cam1) -- (workers);
    \draw[-{Latex}] (cam2) -- (workers);
    \draw[-{Latex}] (workers) -- (yolo);
    \draw[-{Latex}] (yolo) -- (track);
    \draw[-{Latex}] (yolo) -- (color);
    \draw[-{Latex}] (yolo) -- (xml);
\end{tikzpicture}
\caption{Arquitectura del Detector Multi-Cámara}
\end{figure}

El sistema sigue una arquitectura modular con los siguientes componentes:

\textbf{Fuentes de Video}:
\begin{itemize}
    \item CAM\_001: Oracle WebSocket (ws://144.22.56.85:5000/ws/CAM\_001)
    \item CAM\_002: Skyline HLS (extracción con Playwright)
\end{itemize}

\textbf{Procesamiento}:
\begin{itemize}
    \item Worker Threads: worker\_oracle() y worker\_skyline() ejecutan en hilos separados
    \item YOLOv8 Model: detect\_vehicles() - detección de objetos
    \item Tracking System: match\_vehicle\_to\_tracked() - seguimiento único
    \item Color Detection: get\_dominant\_color() - análisis HSV
    \item XML Database: save\_detection\_to\_xml() - persistencia
\end{itemize}

\textbf{Interfaces}:
\begin{itemize}
    \item Flask API: /api/vehicles - estadísticas en JSON
    \item WebSocket: /ws/stream - streaming de video procesado
\end{itemize}

\subsubsection{Configuración de Cámaras}

\begin{lstlisting}[caption={Configuración Multi-Cámara}, style=pythonstyle, label=lst:cam-config]
CAMERAS_CONFIG = {
    'CAM_001': {
        'name': 'Oracle Server',
        'type': 'websocket',
        'url': 'ws://144.22.56.85:5000/ws/CAM_001'
    },
    'CAM_002': {
        'name': 'Skyline Cochabamba',
        'type': 'skyline',
        'url': 'https://www.skylinewebcams.com/...'
    }
}

# Estado global por cámara
camera_states = {
    'CAM_001': {
        'raw_frame': None,
        'processed_frame': None,
        'tracked_vehicles': {},
        'next_vehicle_id': 1,
        'status': 'offline'
    },
    'CAM_002': { ... }
}
\end{lstlisting}

\subsubsection{Algoritmo de Tracking}

El sistema implementa tracking único basado en distancia euclidiana entre centroides:

\begin{enumerate}
    \item Calcular centroide del bbox detectado
    \item Comparar con vehículos tracked existentes
    \item Si distancia \textless{} max\_distance: asociar a vehículo existente
    \item Si no: crear nuevo ID de vehículo
    \item Eliminar vehículos no vistos en max\_frames\_missing frames
\end{enumerate}

\begin{lstlisting}[caption={Algoritmo de Tracking}, style=pythonstyle, label=lst:tracking]
def match_vehicle_to_tracked(bbox, tracked_vehicles, max_distance):
    center = calculate_bbox_center(bbox)
    best_match_id = None
    best_distance = float('inf')
    
    for vehicle_id, vehicle_data in tracked_vehicles.items():
        tracked_center = calculate_bbox_center(vehicle_data['bbox'])
        distance = calculate_distance(center, tracked_center)
        
        if distance < max_distance and distance < best_distance:
            best_distance = distance
            best_match_id = vehicle_id
    
    return best_match_id
\end{lstlisting}

\subsubsection{Detección de Color}

Utiliza espacio de color HSV para determinar color dominante del vehículo:

\begin{table}[H]
\centering
\begin{tabular}{|l|l|}
\hline
\textbf{Rango Hue} & \textbf{Color} \\
\hline
\textless 10 o \textgreater 170 & Rojo \\
10-25 & Naranja \\
25-35 & Amarillo \\
35-85 & Verde \\
85-130 & Azul \\
130-160 & Morado \\
S\textless50, V\textgreater200 & Blanco \\
V\textless50 & Negro \\
Otros & Gris \\
\hline
\end{tabular}
\caption{Mapeo HSV a Colores}
\end{table}

\subsubsection{API Flask}

\begin{table}[H]
\centering
\begin{tabular}{|l|l|p{5cm}|}
\hline
\textbf{Endpoint} & \textbf{Método} & \textbf{Descripción} \\
\hline
/ & GET & Página de información del detector \\
\hline
/api/vehicles & GET & Estadísticas (camera\_id param) \\
\hline
/stats & GET & Stats simplificadas (camera\_id param) \\
\hline
/api/reset & POST & Reset contadores \\
\hline
/ws/stream & WebSocket & Stream de video procesado \\
\hline
\end{tabular}
\caption{Endpoints del Detector Python}
\end{table}

\subsection{voice\_server.py}

\textbf{Archivo}: \texttt{deteccion/voice\_server.py} (180 líneas)

\textbf{Responsabilidad}: Reconocimiento de voz en español usando Vosk, comunicación WebSocket con frontend.

\subsubsection{Arquitect ura de Voz}

El sistema de reconocimiento de voz sigue el siguiente flujo:

\begin{enumerate}
    \item \textbf{Micrófono} captura audio a 16kHz mediante sounddevice
    \item \textbf{Audio Queue} (queue.Queue()) almacena buffers de audio
    \item \textbf{Vosk Model} (KaldiRecognizer) procesa el audio
    \item \textbf{Flask-SocketIO} emite eventos:
        \begin{itemize}
            \item voice\_command: Texto reconocido final
            \item voice\_partial: Texto parcial en tiempo real
        \end{itemize}
    \item \textbf{JavaScript Client} en modulo4.blade.php recibe y procesa comandos
\end{enumerate}

\newpage

\section{Vistas Blade}

\subsection{modulo1.blade.php}

\textbf{Archivo}: \texttt{resources/views/modulo1.blade.php} (323 líneas)

\textbf{Componentes principales}:
\begin{itemize}
    \item Selector de cámara (CAM\_001, CAM\_002)
    \item Canvas 800x600px para video en tiempo real
    \item Estadísticas: vehículos actuales, total detectados, FPS, promedio
    \item Gráficos de tipos y colores de vehículos
\end{itemize}

\subsection{modulo2.blade.php}

\textbf{Archivo}: \texttt{resources/views/modulo2.blade.php}

\textbf{Funcionalidades}:
\begin{itemize}
    \item Filtros por tipo, fecha, cámara
    \item Tabla de detecciones con paginación
    \item Exportación a Excel/CSV
    \item Gráficos estadísticos con Chart.js
\end{itemize}

\subsection{modulo3.blade.php}

\textbf{Archivo}: \texttt{resources/views/modulo3.blade.php}

\textbf{Panel de Administración} (solo admins):
\begin{itemize}
    \item Gestión de usuarios
    \item Configuración de notificaciones (thresholds)
    \item Configuración de reportes automáticos
    \item Limpieza programada de XML
    \item Monitor de estado de cámaras
\end{itemize}

\subsection{modulo4.blade.php}

\textbf{Archivo}: \texttt{resources/views/modulo4.blade.php}

\textbf{Gestión de Comandos de Voz}:
\begin{itemize}
    \item CRUD completo de comandos
    \item Panel de pruebas en tiempo real
    \item Integración con SocketIO para reconocimiento
\end{itemize}

\subsection{modulo5.blade.php}

\textbf{Archivo}: \texttt{resources/views/modulo5.blade.php}

\textbf{Gestión de Usuarios} (solo admins):
\begin{itemize}
    \item Crear/eliminar usuarios
    \item Cambiar roles (user/admin)
    \item Tabla de usuarios activos
\end{itemize}

\newpage
